Vision-language models (VLMs) that connect a pretrained vision encoder to a large language model (LLM) have
become the default way to build general-purpose multimodal systems \cite{yin2024mllmsurvey,zhang2024vlmsurvey,awais2025foundation}. As these systems move into screening, moderation and medical triage, a failure mode
inherited from earlier deep networks becomes critical: \textbf{shortcut learning}, the use of a cue that predicts the
label in training but not in deployment \cite{geirhos2020shortcut}. Shortcuts have already produced high-accuracy
models that read hospital markers instead of lung pathology \cite{degrave2021covid} and watermarks instead of
objects \cite{lapuschkin2019clever}. They fail silently on exactly the under-represented groups where errors are
most costly \cite{brown2023shortcut,ongly2024shortcut}.

A substantial body of work on spurious correlations establishes how to measure and mitigate such failures in
classifiers:
- worst-group accuracy over groups defined by label and spurious attribute \cite{sagawa2020groupdro};
- group-aware or label-free reweighting \cite{liu2021jtt,idrissi2022balancing};
- the finding that features usually retain the core signal, so retraining only the readout recovers most
robustness \cite{izmailov2022feature,kirichenko2023dfr}.

For VLMs, recent studies show that multimodal LLMs inherit spurious cues and fail on minority groups
\cite{hosseini2025spurlens,yang2025spuriverse,zhang2024oodmllm}. They attribute this to the vision encoder
\cite{tong2024eyes}, to data rather than scale \cite{sarridis2026scaling}, or remove it post hoc by unlearning
\cite{sarker2026purge}. In parallel, a separate line of work compares the connectors that map visual features into
the LLM, such as MLP projectors, Q-Formers and gated cross-attention \cite{liu2024llava15,li2023blip2,alayrac2022flamingo}, and ranks them by benchmark accuracy and efficiency \cite{mckinzie2024mm1,cha2024honeybee,karamcheti2024prismatic}.

However, these two lines of work have not met, which leaves three limitations.
1. \textbf{No controlled test of the connector.} Every robustness study of VLMs treats the connector as fixed, and
every connector study measures average accuracy rather than worst-group behaviour. Whether the component
that decides \emph{what visual information reaches the LLM} also decides \emph{whether the model relies on a
shortcut} is untested.
2. \textbf{Conflicting evidence on gated cross-attention.} Zero-initialised, learnable tanh gates are the standard
way to inject a new modality into a pretrained LLM: Flamingo, IDEFICS, LLaMA-Adapter and Llama 3.2 Vision
all use them \cite{alayrac2022flamingo,zhang2024llamaadapter,grattafiori2024llama3}. Cross-attention
injection performs well with frozen backbones \cite{laurencon2024matters}. Yet learnable injection gates can end
up behaviourally closed \cite{farazi2026gate}, and every attention gate in the released Llama 3.2 Vision
checkpoints is below 0.012. The consequence of a near-shut gate for \emph{what the model relies on} has never
been measured. Work on modality imbalance shows that multimodal networks can greedily ignore a modality
\cite{wu2022greedy,huang2022modality}, but not which component of an LLM-based VLM causes it, nor what it does
to spurious-group robustness.
3. \textbf{Results from too few seeds.} Comparisons rest on one or a few training runs, although seed variance
alone can reverse conclusions in deep learning \cite{henderson2018matters,bouthillier2021variance}.

The field therefore cannot say whether a widely used architectural choice silently causes shortcut learning,
or how to avoid it.

To address these limitations, this study tests whether the vision-language connector, and specifically its
injection gate, determines shortcut reliance in LLM-based VLMs.
- \textbf{What is held fixed:} the LLM (Qwen3.5-4B, replicated with Phi-4-mini), frozen SigLIP features, data,
prompts and training budget.
- \textbf{What varies:} only the connector, at four levels: an MLP projector, a Q-Former, and gated
cross-attention with a learnable or a fixed gate. The two gate variants have identical parameters; a third,
learnable gate initialised open separates trainability from initialisation.
- \textbf{Where it is measured:} worst-group accuracy on Waterbirds, a controlled watermark task and CelebA.
- \textbf{How it is tested:} hypotheses formed on a 3-seed pilot and confirmed on up to 12 new paired
seeds (15 in total), with effect sizes,
robust and non-parametric tests, and an equivalence test.
- \textbf{Fixes and rival explanations:} we log the gate during training, read the gates of released IDEFICS
checkpoints, test a connector-only analogue of deep feature reweighting, and probe rival explanations
(token count, encoder scale, frozen LLM, training length).

This paper makes three contributions.
1. \textbf{A transmission mechanism.} We identify the learnable injection gate as a \emph{transmission bottleneck}
between features and readout. Initialised at zero, the gate never opens, and the model then relies on the shortcut; fixing it raises
worst-group accuracy by 10.6 and 25.2 points (d\_z = 1.10 and 1.34). This extends the feature--readout account
of shortcut learning and links gate closure to robustness.
2. \textbf{Variance as an outcome.} The learnable gate turns robustness into a seed lottery, with up to 10 of 15
runs collapsing. Once the gate is fixed, we detect no significant difference between connector types, although equivalence
is not established. Together these establish
that architectural robustness claims require seed-level evaluation.
3. \textbf{A practical fix and its boundary.} Retraining only the connector on a small group-balanced set raises
watermark worst-group accuracy to 89.6\% with near-zero seed variance, but not when the shortcut is entangled
with the object.

Section 2 reviews related work. Section 3 describes the design and statistical protocol. Section 4 reports the
results, Section 5 discusses mechanisms and implications, and Section 6 concludes. Code, per-seed results and
figures are public.
